\section{The mesh-resolution sensitivity envelope}
\label{app:envelope}

The envelope quoted throughout this paper as $0.9$--$5.9\%$ is the range over
which the normalized geopotential dispersion $\potvarn$ moves when the same
body is meshed at different resolutions. It is used as a scale against which
an improvement may be judged large or small, and it is not a statistical
threshold: nothing about it licenses a significance statement. Because
several sections compare an improvement against the value the envelope takes
on one particular body, this appendix records those values, how they were
obtained, and how far they depend on choices that are not forced.

\subsection{Definition and provenance}
\label{app:envelope:def}

Paper~IV decimated four shape models---Eros, Itokawa, Gaspra and Ida---over a
ladder from roughly $2{,}000$ to $42{,}000$ facets and recorded $\potvarn$ and
the area-weighted mean slope at each rung. For each body the envelope is the
full range of $\potvarn$ across its ladder, as a fraction of the smallest
value:
\begin{equation}
E \;=\; \frac{\max_i \potvarn^{(i)} - \min_i \potvarn^{(i)}}
             {\min_i \potvarn^{(i)}} ,
\label{eq:envelope}
\end{equation}
the index running over the rungs. The published figures are $0.9\%$ at one
extreme and $5.9\%$ at the other, and the paper quotes that pair as a range
across the test set rather than as a property of any one body.

Two things follow that matter here. The envelope is a property of a body and
its ladder, not of the method, so comparing an improvement on Itokawa against
the largest value in a four-body set is the weaker of the two available
comparisons; and the ladder itself is a choice, so the envelope moves when the
decimator does.

\subsection{Recomputed values}
\label{app:envelope:values}

Table~\ref{tab:envelope} gives Eq.~\eqref{eq:envelope} recomputed for this
paper on the same four bodies, with both decimators, on ladders of six rungs
each. The dispersion column is the envelope; the slope column is the same
quantity for the mean surface slope, included because the comparison between
the two is the result of Paper~IV; and the plateau column is the change in
$\potvarn$ between the two finest meshes alone.

\begin{table}[htbp]
\centering
\small
\setlength{\tabcolsep}{5pt}
\caption{The envelope of Eq.~\eqref{eq:envelope} recomputed on the four
Paper~IV resolution bodies, six rungs each. The slope column is the same
range for the area-weighted mean slope. The plateau is the change in
$\potvarn$ between the two finest meshes, which is the quantity Paper~IV used
to argue that the dispersion has settled. Values in bold are those quoted
elsewhere in this paper: the Itokawa envelope in
Sections~\ref{sec:itokawa} and~\ref{sec:conclusions}, and the Eros envelope
in Sections~\ref{sec:nulls} and~\ref{sec:failures}.}
\label{tab:envelope}
\begin{tabular}{llrrrr}
\toprule
Decimator & Body & Facets & Dispersion $E$ & Slope range & Plateau \\
\midrule
Vertex clustering & Eros    & $2{,}097$--$42{,}907$ & $5.40\%$ & $21.91\%$ & $0.083\%$ \\
                  & Itokawa & $2{,}183$--$48{,}785$ & $0.95\%$ & $14.24\%$ & $0.019\%$ \\
                  & Gaspra  & $2{,}821$--$32{,}040$ & $1.48\%$ & $1.84\%$  & $0.000\%$ \\
                  & Ida     & $2{,}141$--$32{,}040$ & $4.40\%$ & $5.79\%$  & $0.009\%$ \\
\addlinespace
Quadric collapse  & Eros    & $2{,}000$--$42{,}000$ & $\mathbf{6.03\%}$ & $3.50\%$ & $0.098\%$ \\
                  & Itokawa & $2{,}000$--$42{,}000$ & $\mathbf{1.35\%}$ & $4.50\%$ & $0.008\%$ \\
                  & Gaspra  & $2{,}000$--$32{,}040$ & $2.92\%$ & $3.15\%$  & $0.000\%$ \\
                  & Ida     & $2{,}000$--$32{,}040$ & $1.91\%$ & $6.19\%$  & $0.000\%$ \\
\bottomrule
\end{tabular}
\end{table}

\subsection{What the table shows}
\label{app:envelope:reading}

\emph{The published range is reproduced at its lower end and not at its
upper.} With vertex clustering, which Paper~IV used, the four dispersion
envelopes run from $0.95\%$ to $5.40\%$ against the published
$0.9$--$5.9\%$. The lower end agrees; the upper is $0.5$ percentage points
short, and the difference is the ladder, whose rungs are set by the
decimator's own grid and do not fall at the same facet counts as Paper~IV's.
We quote $0.9$--$5.9\%$ throughout because it is the published figure, and
note here that the recomputed range is slightly narrower.

\emph{The envelope depends on the decimator, and not weakly.} On Eros it is
$5.40\%$ with clustering and $6.03\%$ with quadric collapse; on Itokawa,
$0.95\%$ and $1.35\%$. That dependence is the reason
Sections~\ref{sec:nulls:misplaced} and~\ref{sec:failures:plane} state both
values when the envelope is used as a test: the misplaced Eros plane returns
an improvement of $5.526\%$, which is $0.92$ times the quadric envelope and
rejected by it, and $1.02$ times the clustering envelope and not rejected.
A diagnostic whose verdict turns on a choice of decimation is worth less than
its arithmetic suggests, and we rely on plane provenance instead.

\emph{The slope column behaves differently under the two decimations.} With
clustering the mean slope varies by $1.8\%$ to $21.9\%$ across the four
bodies and rises monotonically with facet count on every one; with quadric
collapse it varies by $3.1\%$ to $6.2\%$ and shows no monotone trend. The
analytic ellipsoid of Section~\ref{sec:discussion:ellipsoid} finds the two
decimations equally accurate on a smooth surface, so the monotone drift under
clustering is a property of rough surfaces that the ellipsoid cannot
reproduce, and we make no claim about which decimation is closer to the truth
on a real body. What survives either way is the comparison Paper~IV drew: the
dispersion moves by a few per cent across a ladder that moves the slope by
several times as much, and the ellipsoid puts that ratio at $19\pm4$ to
$59\pm18$ depending on resolution.

\emph{The plateau is tighter than the figure the paper uses.} Between the two
finest meshes the dispersion moves by $0.008\%$ on Itokawa and by at most
$0.098\%$ on any of the four bodies, with either decimator. Paper~IV reports
$0.60\%$ for this quantity, and the curvature sensitivity intervals of
Section~\ref{sec:itokawa} are evaluated at that level. Retaining the
published figure makes those intervals conservative by roughly a factor of
six, which is the direction in which an interval should err.

\subsection{Scope}
\label{app:envelope:scope}

Four bodies, six rungs, two decimators. The envelope is not known for Bennu,
67P or Arrokoth, whose meshes were decimated once each rather than laddered,
so the comparisons in Section~\ref{sec:nulls} for those three bodies use the
cross-body range and not a body-specific value. Extending the ladder to them
would cost one batch of runs and would sharpen three of the four
weak-signal cases; we did not do it, and a reader who wants a body-specific
envelope for those bodies should treat its absence as a gap rather than as
evidence that the cross-body range applies.
